Synthetic speech, singing, music and environmental sound can now be generated well enough to fool human listeners. A detector that is useful outside a benchmark has to cope with three things at once: it does not know the audio type in advance, the recording may have been compressed, resampled or filtered on its way to the detector, and it may have to run on a phone or another small device. Large self-supervised models score well on benchmarks but are too heavy for that last requirement, so we look at a small network whose robustness is built in during training and costs nothing at inference.

The study evaluates audio deepfake detection on the multi-type AT-ADD Track~2 benchmark. In clean baseline selection across three seeds, AudioSet-pretrained PANNs MobileNetV2 (E3, 3.54\,M parameters) achieves the highest clean performance (type-averaged Macro-F1 0.627, AUROC 0.880, EER 0.199) and lowest variance over compact scratch candidates, and is selected as reference baseline M0. In exploratory regularisation experiments on a compact MobileNetV3-Small network (about 1.52\,M parameters), we compare a clean-trained detector with Frequency-Time Robustness Coherence Learning (FT-RCL), which trains the network on pairs of clean and degraded clips with prediction and feature-signature agreement losses, and with an extension, STAR-Edge. All regularisation models were trained on a stratified subset of 13{,}800 training clips and evaluated on 13{,}287 evaluation clips drawn from all 38 generators that never appear in training.

On the official type-averaged Macro-F1, the baseline scores 0.587 (95\% bootstrap interval 0.577--0.598), FT-RCL 0.589 and STAR-Edge 0.599. FT-RCL does not improve on the baseline: its pooled Macro-F1 is 0.013 lower and its EER 1.8 points higher, and the paired bootstrap intervals of both differences exclude zero. STAR-Edge gives the best AUROC (0.866 against 0.852) at an EER equal to the baseline. All three models are close to chance on music and sound. The gap between development scores (0.91--0.93) and evaluation scores (0.66--0.67) is explained by the benchmark itself: the development split shares every fake generator with the training split, whereas the evaluation split shares none.

We also audited our own pipeline. An exploratory pilot robustness test was evaluated on fake music clips, and the feature-signature module constrained only three of its eight bands. No latency or quantised-accuracy measurements were made on physical edge hardware; the compact MobileNetV3-Small network requires 36.7 million multiply-accumulates per clip and 6.08\,MB of FP32 weights.
